% ── SECTION 1: INTRODUCTION — THE TWO-WINGS TEST (~700 words) ───────────────

\section{Introduction: The Two-Wings Test}
\label{sec:intro}

% Key question: what would genuine confirmation of the two-wings principle
% look like? Not shared conclusions — but shared structure.
% Introduce structural isomorphism vs. analogy.
% Preview six convergences and chronological sequence.

\AbdulBaha taught that science and religion are ``like the two wings upon which
the bird of human intelligence can soar into the heaven of truth.''
The principle is foundational to Bahá'í teaching and frequently invoked.
But what would it mean to confirm it?

Previous Bahá'í-science dialogue has often worked by identifying resonances:
quantum uncertainty sounds like the unknowability of God; entanglement sounds
like the interconnection of all things; the Big Bang sounds like creation.
These resonances are suggestive.
But they are analogical, not structural.
Analogy shows that two things look similar; structural isomorphism shows that
they have the same logical form.

This paper asks whether the two-wings principle can be confirmed at the level
of structure --- whether the Bahá'í metaphysical framework and the framework
that modern physics has empirically arrived at share the same logical
architecture, arrived at by independent means.

The answer we document is yes.
The Bahá'í metaphysics of \mahabbat --- relational love as the creative
foundation of existence at every degree --- is the same logical structure
as relational quantum mechanics~\citep{Rovelli1996} and process
ontology~\citep{Whitehead1929}: existence constituted by relationship,
properties emerging through relational interaction, the force sustaining
relationship identified as love.

% TODO: preview the six convergences, the chronological sequence,
% and the distinction between confirmation and appropriation.
